% 车轮效应教学构造：8根辐条，每帧逆时针转42度，辐条间隔45度。宽约107mm。
\begin{tikzpicture}[x=1mm,y=1mm,font=\figurefont\small,
  rim/.style={draw=FigureStroke,line width=.8pt},
  spoke/.style={draw=FigureGrid,line width=.7pt},
  lab/.style={font=\figurefont\footnotesize,text=FigureMuted}]
  \foreach \k/\cx in {0/0,1/28,2/56}{
    \begin{scope}[shift={(\cx,0)}]
      \draw[FigureMuted,dashed,line width=.6pt] (0,0) -- (0,14);
      \foreach \s in {0,...,7}{
        \draw[spoke] (0,0) -- ({90+42*\k+45*\s}:10.5);
      }
      \draw[rim] (0,0) circle (10.5);
      \ifnum\k>0
        \draw[FigureBlue,line width=1.5pt] (0,0) -- ({90-3*\k}:10.5);
      \fi
      \draw[FigureRed,line width=1.5pt] (0,0) -- ({90+42*\k}:10.5);
      \fill[FigureStroke] (0,0) circle (.8);
      \node at (0,-14.5) {第\k 帧};
    \end{scope}
  }
  \draw[-{Stealth[length=1.6mm]},FigureRed,line width=.6pt]
    (28,0) ++(90:6) arc[start angle=90,end angle=132,radius=6];
  \node[lab,anchor=east] at (17.5,11) {$+42^\circ$};
  \node[lab,anchor=west] at (29.5,13) {$-3^\circ$};
  \node[lab,anchor=west] at (57.5,13) {$-6^\circ$};
  \draw[FigureRed,line width=1.5pt] (70,6) -- (75,6);
  \node[anchor=west,align=left] at (76,6) {同一根辐条};
  \draw[FigureBlue,line width=1.5pt] (70,-1) -- (75,-1);
  \node[anchor=west,align=left] at (76,-1) {最接近参考线\\的辐条};
\end{tikzpicture}
